% Method exclusion only; no new main-gate closure or priority claim.
\subsubsection*{An integral obstruction for the proposed Lagrangian-surface method}

Let $(A,L)$ be a complex abelian surface with
$\operatorname{End}(A)=\mathbb Z$ and
$\operatorname{NS}(A)=\mathbb Z[L]$, where $L$ has type $(1,2)$.
Put $\lambda=\lambda_L$ and let
$\lambda_d:\widehat A\to A$ be the primitive dual polarization,
so that $\lambda_d\lambda=[2]_A$ and
$\lambda\lambda_d=[2]_{\widehat A}$.
Suppose that an abelian variety $B$ admits an isogeny
\[
                     \pi:A^n\longrightarrow B
\]
of odd degree. For every integral curve cycle $z$ on $B$ and
every homomorphism $f:B\to A\times\widehat A$, the mixed
endomorphism $A\to A$ of the symmetric curve homomorphism
$\Phi_{f_*z}:\widehat A\times A\to A\times\widehat A$
is an even integer multiple of the identity. The same statement
holds for an arbitrary morphism, since translation does not change
this homomorphism. In particular, it applies to every divisor
cycle on a smooth projective surface whose Albanese variety is
such a $B$. We make no assumption that the divisor classes of
the surface come from its ambient abelian variety.

\begin{proof}
For a curve normalization $C\to V$ in an abelian variety, let
$u:J(C)\to V$ be the induced homomorphism. We use the integral
symmetric homomorphism
\[
             \Phi_C=u\lambda_{J(C)}^{-1}\widehat u
                         :\widehat V\longrightarrow V,
\]
and extend it additively to signed curve cycles. These are the
Matsusaka homomorphisms already used in C.4r. In particular,
$\Phi_{g_*z}=g\Phi_z\widehat g$, and translation acts trivially.

The primitive Hom-lattice calculation in C.4t gives
\[
 \operatorname{Hom}(A,\widehat A)=\mathbb Z\lambda,
 \qquad
 \operatorname{Hom}(\widehat A,A)=\mathbb Z\lambda_d.
\]
For a curve cycle on $A^n$ we can therefore write
\[
                         \Phi_z=\lambda_d S,
             \qquad S\in\operatorname{Mat}_n(\mathbb Z),
             \qquad S^t=S.
\]
Every homomorphism $g:A^n\to A\times\widehat A$ has the form
\[
                  g=\begin{pmatrix}r\\ \lambda s\end{pmatrix},
                     \qquad r,s\in\mathbb Z^{1\times n}.
\]
The upper-right block of $g\Phi_z\widehat g$ is
\[
               r\lambda_d S\lambda s^t
                       =2rSs^t\,\operatorname{id}_A.
\]
It is thus even for every integral curve cycle on $A^n$.

Write $d=\deg\pi$, and pull a curve cycle $z$ on $B$ back by
the finite flat isogeny $\pi$. The identity
$\pi_*\pi^*z=dz$ is integral. Applying the preceding calculation
to $g=f\pi$ shows that $d$ times the mixed integer of
$\Phi_{f_*z}$ is even. Since $d$ is odd, that integer is even.
An odd-degree isogeny in the opposite direction would give the
same hypothesis: an odd integer killing its kernel gives an
integral reverse isogeny of odd degree.

Finally, a morphism from a smooth projective surface to an
abelian variety factors through its Albanese map, up to
translation. Pushing each irreducible divisor curve through
this factorization proves the assertion for all divisor cycles.
\end{proof}

This excludes two specific full-span surface constructions as a
source of the missing odd Poincare curve coefficient.
\citet[Theorem 6.8, author pp.~25--26]{frC4BastianelliPirolaStoppino2010}
prove that their Lagrangian Galois-closure surface associated
with a type-$(1,2)$ surface $A$ has Albanese variety $A^2$.
The lemma therefore excludes every divisor on that surface,
including divisor classes not inherited from $A^2$.

For the newer construction of
\citet[Proposition 4.2, Theorem 5.6, Proposition 6.2, and
Corollary 6.3]{frC4GrossiMoretti2025}, let $A_0$ have type
$(1,6)$ and let $K\simeq(\mathbb Z/3)^2$ be their subgroup.
For a subgroup $G\subset K$ of order three, their quotient
$A=A_0/G$ has type $(1,2)$, while the Albanese variety of
their quotient surface is
\[
                  B=(A_0\times A_0)/\operatorname{diag}K.
\]
Assume that this resulting $A$ has the generic End and
Neron--Severi groups used above. The map
$q^2:A_0^2\to A^2$ has odd degree and its kernel is killed
by three. Consequently $[3]_{A_0^2}$ factors through an
odd-degree isogeny $u:A^2\to A_0^2$. Composing $u$ with
the displayed quotient gives an odd-degree isogeny
$A^2\to B$. Every divisor on their quotient surface is
therefore covered by the same obstruction.

The required numerical Poincare curve class of C.4r has
mixed coefficient $\pm1$, so these constructions cannot supply
it. This is a restriction on this particular method. It does
not show that arbitrary curves on $A\times\widehat A$ have
even mixed coefficient, and it does not refute or prove the
remaining integral Jacobian-product splitting assertion.
